% --- LaTeX Homework Template - S. Venkatraman ---

% --- Set document class and font size ---

\documentclass[letterpaper, 11pt]{article}

% --- Package imports ---

\usepackage{
  amsmath, amsthm, amssymb, mathtools, dsfont,	  % Math typesetting
  graphicx, wrapfig, subfig, float,                  % Figures and graphics formatting
  listings, color, inconsolata, pythonhighlight,     % Code formatting
  fancyhdr, sectsty, hyperref, enumerate, enumitem } % Headers/footers, section fonts, links, lists

% --- Page layout settings ---

% Set page margins
\usepackage[left=1.35in, right=1.35in, bottom=1in, top=1.1in, headsep=0.2in]{geometry}

% Anchor footnotes to the bottom of the page
\usepackage[bottom]{footmisc}

% Set line spacing
\renewcommand{\baselinestretch}{1.2}

% Set spacing between paragraphs
\setlength{\parskip}{1.5mm}

% Allow multi-line equations to break onto the next page
\allowdisplaybreaks

% Enumerated lists: make numbers flush left, with parentheses around them
\setlist[enumerate]{wide=0pt, leftmargin=21pt, labelwidth=0pt, align=left}
\setenumerate[1]{label={(\arabic*)}}

% --- Page formatting settings ---

% Set link colors for labeled items (blue) and citations (red)
\hypersetup{colorlinks=true, linkcolor=blue, citecolor=red}

% Make reference section title font smaller
\renewcommand{\refname}{\large\bf{References}}

% --- Settings for printing computer code ---

% Define colors for green text (comments), grey text (line numbers),
% and green frame around code
\definecolor{greenText}{rgb}{0.5, 0.7, 0.5}
\definecolor{greyText}{rgb}{0.5, 0.5, 0.5}
\definecolor{codeFrame}{rgb}{0.5, 0.7, 0.5}

% Define code settings
\lstdefinestyle{code} {
  frame=single, rulecolor=\color{codeFrame},            % Include a green frame around the code
  numbers=left,                                         % Include line numbers
  numbersep=8pt,                                        % Add space between line numbers and frame
  numberstyle=\tiny\color{greyText},                    % Line number font size (tiny) and color (grey)
  commentstyle=\color{greenText},                       % Put comments in green text
  basicstyle=\linespread{1.1}\ttfamily\footnotesize,    % Set code line spacing
  keywordstyle=\ttfamily\footnotesize,                  % No special formatting for keywords
  showstringspaces=false,                               % No marks for spaces
  xleftmargin=1.95em,                                   % Align code frame with main text
  framexleftmargin=1.6em,                               % Extend frame left margin to include line numbers
  breaklines=true,                                      % Wrap long lines of code
  postbreak=\mbox{\textcolor{greenText}{$\hookrightarrow$}\space} % Mark wrapped lines with an arrow
}

% Set all code listings to be styled with the above settings
\lstset{style=code}

% --- Math/Statistics commands ---

% Add a reference number to a single line of a multi-line equation
% Usage: "\numberthis\label{labelNameHere}" in an align or gather environment
\newcommand\numberthis{\addtocounter{equation}{1}\tag{\theequation}}

% Shortcut for bold text in math mode, e.g. $\b{X}$
\let\b\mathbf

% Shortcut for bold Greek letters, e.g. $\bg{\beta}$
\let\bg\boldsymbol

% Shortcut for calligraphic script, e.g. %\mc{M}$
\let\mc\mathcal

% \mathscr{(letter here)} is sometimes used to denote vector spaces
\usepackage[mathscr]{euscript}

% Convergence: right arrow with optional text on top
% E.g. $\converge[w]$ for weak convergence
\newcommand{\converge}[1][]{\xrightarrow{#1}}

% Normal distribution: arguments are the mean and variance
% E.g. $\normal{\mu}{\sigma}$
\newcommand{\normal}[2]{\mathcal{N}\left(#1,#2\right)}

% Uniform distribution: arguments are the left and right endpoints
% E.g. $\unif{0}{1}$
\newcommand{\unif}[2]{\text{Uniform}(#1,#2)}

% Independent and identically distributed random variables
% E.g. $ X_1,...,X_n \iid \normal{0}{1}$
\newcommand{\iid}{\stackrel{\smash{\text{iid}}}{\sim}}

% Equality: equals sign with optional text on top
% E.g. $X \equals[d] Y$ for equality in distribution
\newcommand{\equals}[1][]{\stackrel{\smash{#1}}{=}}

% Math mode symbols for common sets and spaces. Example usage: $\R$
\newcommand{\R}{\mathbb{R}}   % Real numbers
\newcommand{\C}{\mathbb{C}}   % Complex numbers
\newcommand{\Q}{\mathbb{Q}}   % Rational numbers
\newcommand{\Z}{\mathbb{Z}}   % Integers
\newcommand{\N}{\mathbb{N}}   % Natural numbers
\newcommand{\F}{\mathcal{F}}  % Calligraphic F for a sigma algebra
\newcommand{\El}{\mathcal{L}} % Calligraphic L, e.g. for L^p spaces

% Math mode symbols for probability
\newcommand{\pr}{\mathbb{P}}    % Probability measure
\newcommand{\E}{\mathbb{E}}     % Expectation, e.g. $\E(X)$
\newcommand{\var}{\text{Var}}   % Variance, e.g. $\var(X)$
\newcommand{\cov}{\text{Cov}}   % Covariance, e.g. $\cov(X,Y)$
\newcommand{\corr}{\text{Corr}} % Correlation, e.g. $\corr(X,Y)$
\newcommand{\B}{\mathcal{B}}    % Borel sigma-algebra

% Other miscellaneous symbols
\newcommand{\tth}{\text{th}}	% Non-italicized 'th', e.g. $n^\tth$
\newcommand{\Oh}{\mathcal{O}}	% Big-O notation, e.g. $\O(n)$
\newcommand{\1}{\mathds{1}}	% Indicator function, e.g. $\1_A$

% Additional commands for math mode
\DeclareMathOperator*{\argmax}{argmax}    % Argmax, e.g. $\argmax_{x\in[0,1]} f(x)$
\DeclareMathOperator*{\argmin}{argmin}    % Argmin, e.g. $\argmin_{x\in[0,1]} f(x)$
\DeclareMathOperator*{\spann}{Span}       % Span, e.g. $\spann\{X_1,...,X_n\}$
\DeclareMathOperator*{\bias}{Bias}        % Bias, e.g. $\bias(\hat\theta)$
\DeclareMathOperator*{\ran}{ran}          % Range of an operator, e.g. $\ran(T) 
\DeclareMathOperator*{\dv}{d\!}           % Non-italicized 'with respect to', e.g. $\int f(x) \dv x$
\DeclareMathOperator*{\diag}{diag}        % Diagonal of a matrix, e.g. $\diag(M)$
\DeclareMathOperator*{\trace}{trace}      % Trace of a matrix, e.g. $\trace(M)$

% Numbered theorem, lemma, etc. settings - e.g., a definition, lemma, and theorem appearing in that 
% order in Section 2 will be numbered Definition 2.1, Lemma 2.2, Theorem 2.3. 
% Example usage: \begin{theorem}[Name of theorem] Theorem statement \end{theorem}
\theoremstyle{definition}
\newtheorem{theorem}{Theorem}[section]
\newtheorem{proposition}[theorem]{Proposition}
\newtheorem{lemma}[theorem]{Lemma}
\newtheorem{corollary}[theorem]{Corollary}
\newtheorem{definition}[theorem]{Definition}
\newtheorem{example}[theorem]{Example}
\newtheorem{remark}[theorem]{Remark}

% Un-numbered theorem, lemma, etc. settings
% Example usage: \begin{lemma*}[Name of lemma] Lemma statement \end{lemma*}
\newtheorem*{theorem*}{Theorem}
\newtheorem*{proposition*}{Proposition}
\newtheorem*{lemma*}{Lemma}
\newtheorem*{corollary*}{Corollary}
\newtheorem*{definition*}{Definition}
\newtheorem*{example*}{Example}
\newtheorem*{remark*}{Remark}
\newtheorem*{claim}{Claim}

% --- Left/right header text (to appear on every page) ---

% Include a line underneath the header, no footer line
\pagestyle{fancy}
\renewcommand{\footrulewidth}{0pt}
\renewcommand{\headrulewidth}{0.4pt}

% Left header text: course name/assignment number
\lhead{MATH-SHU 238 (HTOP) -- Homework 4}

% Right header text: your name
\rhead{Yixia Yu (yy5091@nyu.edu)}

% --- Document starts here ---

\begin{document}
\subsection*{Problem 1}
Uniform distribution. A random variable that is equally likely to take any value in a finite set $S$ is said to have the uniform distribution on $S$. If $U$ is such a random variable and $\emptyset \neq R \subseteq S$, show that the distribution of $U$ conditional on $\{U \in R\}$ is uniform on $R$.
\begin{proof}
    Since $U$ is uniformly distributed on $S$, we have:
\end{proof}
\subsection*{Problem 2}
\begin{align*}
    & 3.\quad\mathrm{Let~}X\text{ be a random variable with distribution function} \\
    & \mathbb{P}\{X\leq x\}=\left\{
   \begin{array}
   {ll}0 & \mathrm{if~}x\leq0, \\
   x & \mathrm{if~}0<x\leq1, \\
   1 & \mathrm{if~}x>1.
   \end{array}\right. \\
    & \mathrm{Let~}F\text{ be a distribution function which is continuous and strictly increasing. Show that }Y=F^{-1}(X) \\
    & \text{is a random variable having distribution function }F.\\
    & \text{ Is it necessary that }F\text{ be continuous and/or strictly}\text{ increasing?}
   \end{align*}
   \begin{proof}
    
   \end{proof}
\subsection*{Problem 3}
(a) Show that any discrete random variable may be written as a linear combination of indicator variables.
\\(b) Show that any random variable may be expressed as the limit of an increasing sequence of discrete random variables.
\\(c) Show that the limit of any increasing convergent sequence of random variables is a random variable.
\begin{proof}
    \end{proof}
\subsection*{Problem 4}
Express the distribution functions of

$$X^{+} = \max\{0, X\}, \quad X^{-} = -\min\{0, X\}, \quad |X| = X^{+} + X^{-}, \quad -X,$$

in terms of the distribution function $F$ of the random variable $X$.
\begin{proof}
    \end{proof}
\subsection*{Problem 5}
The real number $m$ is called a median of the distribution function $F$ whenever $\lim_{y \to m} F(y) \leq \frac{1}{2} \leq F(m)$.

(a) Show that every distribution function $F$ has at least one median, and that the set of medians of $F$ is a closed interval of $\mathbb{R}$.

(b) Show, if $F$ is continuous, that $F(m)=\frac{1}{2}$ for any median $m$.
\begin{proof}
    \end{proof}
\subsection*{Problem 6}
Let $(\Omega, \mathcal{F})$ be a measurable space, $A_{1}$,$\cdots$$A_{n} \in \mathcal{F}$ and $b_{1}$,$\cdots$$b_{n} \in \mathbb{R}$. Set $X = b_{1}1_{A_{1}} + \cdots + b_{n}1_{A_{n}}$. Suppose that for all $i \neq j$, $A_{i} \cap A_{j} = \emptyset$ and $b_{i} \neq b_{j}$. Show that $\sigma(X) = \sigma(\{A_{1}, \cdots, A_{n}\})$.
\begin{proof}
    \end{proof}
\subsection*{Problem 7}
Consider the random walk on $\{0,1,\cdotp\cdotp\cdotp N\}$ that reflects at 0, introduced in the class. More precisely, let $S_0=k,S_{j+1}=S_j+X_j$, where $X_j=1$ if $S_j=0$, otherwise $\mathbb{P}[X_j=1]=p$ and $\mathbb{P}[X_j=-1]=q.$ Compute the expected number of steps to frst reach state $N$, as a $\tilde{\text{function of }k}.\tilde{\text{ Distinguish the cases }p}=q$ and $p\neq q.$
\begin{proof}
    \end{proof}
\subsection*{Problem 8}
If $F: \mathbb{R} \to [ 0, 1]$ satisfes monotone increasing, right continuity, $\lim _{x\to - \infty }F( x) = 0$ and $\lim_{x\to+\infty}F(x)=1$, show that $F$ is the distribution function of some random variable on some probability space. (Hint: vou may want to revisit Section 2.3, excercise 3)
\begin{proof}
    \end{proof}

\end{document}
